\newpage
\section{Trabajo Extra / Voluntario: ILS-ES con Reinicialización por Catástrofe (CHC/GADEGD)}
\label{sec:extra}

\subsection{Motivación: Convergencia Prematura en ILS-ES}

Todo algoritmo de búsqueda metaheurística se enfrenta al dilema fundamental de equilibrar la
\textbf{exploración} —búsqueda en amplitud hacia regiones inexploradas— y la \textbf{explotación}
—profundización en regiones prometedoras ya identificadas—. En el ILS-ES estándar presentado en la
Sección~\ref{sec:algoritmos}, el criterio de aceptación empleado es el del \emph{mejor}: únicamente se
acepta una nueva solución si supera estrictamente a la mejor histórica.

Este criterio favorece una explotación muy agresiva que resulta eficaz en las primeras iteraciones,
pero que a largo plazo provoca la \textbf{convergencia prematura}: la búsqueda acaba concentrada en
la cuenca de atracción del primer óptimo local de calidad aceptable que halla, sin capacidad de
escapar hacia regiones del espacio de soluciones que podrían ofrecer mejores óptimos. La diversidad
se extingue de forma irreversible, y el algoritmo queda bloqueado.

\subsection{Fundamento Teórico: CHC y GADEGD}

La literatura en algoritmos evolutivos propone distintos mecanismos para inyectar diversidad de
forma controlada cuando se detecta estancamiento. En este trabajo se ha incorporado la inspiración
de dos enfoques clásicos:

\subsubsection{Algoritmo CHC (Eshelman \& Schaffer, 1993)}

El algoritmo \textbf{CHC} (\textit{Cross generational elitist selection, Heterogeneous recombination,
and Cataclysmic mutation}) introduce el concepto de \emph{divergencia cataclísmica}: cuando el
sistema detecta que ninguna solución de la población mejora durante un número determinado de
generaciones, aplica una \textbf{macromutación extrema} sobre la mejor solución histórica, alterando
entre el 35\% y el 50\% de sus componentes. Esta perturbación masiva fuerza al algoritmo a saltar a
una cuenca de atracción completamente distinta y retomar la exploración desde una posición
geográficamente alejada del óptimo local previo.

El mecanismo CHC se estructura en cuatro fases:
\begin{enumerate}[noitemsep]
  \item \textbf{Monitorización del estancamiento:} se cuenta el número de iteraciones consecutivas
  sin mejora de la mejor solución global.
  \item \textbf{Detección del umbral:} si el contador supera el límite $\texttt{it}_0$
  (configurable; valor por defecto 10), se considera que el algoritmo está estancado.
  \item \textbf{Fase de divergencia (catástrofe):} en lugar de aplicar la mutación estándar del
  20\%, se aplica una macromutación del 40\% sobre el campeón histórico.
  \item \textbf{Reinicio del contador:} si desde la nueva región se logra mejora, el contador se
  resetea a cero y se retoma la intensificación normal.
\end{enumerate}

\subsubsection{Refinamiento GADEGD (Diversificación Voraz)}

El modelo \textbf{GADEGD} va un paso más allá y propone que la reinicialización no se realice a
ciegas. En lugar de aplicar una macromutación aleatoria, el algoritmo inyecta una solución nueva
construida mediante un \textbf{procedimiento Greedy Aleatorizado}: se seleccionan activos con
preferencia por los de mayor rentabilidad histórica, añadiendo aleatoriedad para inducir diversidad.
Así, la reinicialización no destruye el conocimiento heurístico acumulado, sino que lo combina con
exploración.

En la implementación de esta práctica se ha optado por la \textbf{Opción A (CHC)}: macromutación
sobre el campeón histórico con tasa del 40\%. La Opción B (GADEGD) queda disponible en el código
fuente pero desactivada por defecto, ya que la Opción A presenta mejor equilibrio entre coste
computacional y calidad en los mercados evaluados.

\subsection{Diseño e Integración en ILS-ES}

\subsubsection{Componentes Adicionales}

La variante ILS-ES+CHC extiende el ILS-ES estándar con los siguientes parámetros adicionales:

\begin{itemize}[noitemsep]
  \item \textbf{\texttt{limit\_stagnation}}: número máximo de iteraciones consecutivas sin mejora
  antes de activar la divergencia. Valor: 10.
  \item \textbf{Tasa de mutación estándar}: 20\% (idéntica al ILS-ES base).
  \item \textbf{Tasa de macromutación}: 40\% (divergencia CHC).
\end{itemize}

\subsubsection{Pseudocódigo}

\begin{algorithm}[H]
\SetAlgoLined
\KwIn{Problema $P$, parámetros ES, $\texttt{limit\_stagnation}=10$}
\KwOut{Mejor solución encontrada}

$S \leftarrow$ Generar-Solución-Aleatoria()\;
$Best \leftarrow$ Enfriamiento-Simulado($S$, parámetros\_ES)\;
$ctr\_estanc \leftarrow 0$\;

\While{quepan evaluaciones}{
  \eIf{$ctr\_estanc \geq \texttt{limit\_stagnation}$}{
    \tcp{Fase de divergencia (CHC): macromutación del 40\%}
    $S' \leftarrow$ Macromutar($Best$, $0.40$)\;
    $ctr\_estanc \leftarrow 0$\;
  }{
    \tcp{Fase de intensificación (ILS estándar): mutación del 20\%}
    $S' \leftarrow$ Mutar($Best$, $0.20$)\;
  }
  Reparar($S'$)\;
  $S'' \leftarrow$ Enfriamiento-Simulado($S'$, parámetros\_ES)\;
  \eIf{$f(S'') > f(Best)$}{
    $Best \leftarrow S''$\;
    $ctr\_estanc \leftarrow 0$\;
  }{
    $ctr\_estanc \leftarrow ctr\_estanc + 1$\;
  }
}
\Return $Best$\;
\caption{ILS-ES con Reinicialización por Catástrofe (CHC/GADEGD)}
\label{alg:ilses_chc}
\end{algorithm}

La diferencia clave respecto al ILS-ES estándar (Algoritmo~\ref{alg:ils-es}) radica en el bloque
condicional que, en función del contador de estancamiento, elige entre aplicar la perturbación
habitual o la macromutación de catástrofe. El operador Reparar garantiza la factibilidad de la
solución mutada, al igual que en el ILS-ES base.

\subsection{Resultados Experimentales}

A continuación se presentan los resultados comparativos entre el ILS-ES estándar y el
ILS-ES+CHC, obtenidos bajo las mismas condiciones experimentales de la Sección~\ref{sec:experimentos}
(50 ejecuciones independientes, semilla base 42, hasta 10\,000 evaluaciones).

\subsubsection*{IBEX 35 ($n=30$)}

\begin{table}[H]
\centering
\caption{Comparación ILS-ES vs.\ ILS-ES+CHC -- IBEX 35}
\label{tab:results_ilses_chc_ibex}
\resizebox{\textwidth}{!}{%
\begin{tabular}{lccccc}
\hline
Algoritmo & Fitness(2015--2024) & Fitness(2025) & Std & Evaluaciones & Tiempo (s) \\
\hline
ILS-ES      & -4.649 & -3.127 & 0.014 & 3130 & 0.003 \\
ILS-ES+CHC  & -4.647 & -3.129 & 0.008 & 8591 & 0.008 \\
\hline
\end{tabular}
}
\end{table}

\subsubsection*{S\&P 100 ($n=97$)}

\begin{table}[H]
\centering
\caption{Comparación ILS-ES vs.\ ILS-ES+CHC -- S\&P 100}
\label{tab:results_ilses_chc_sp100}
\resizebox{\textwidth}{!}{%
\begin{tabular}{lccccc}
\hline
Algoritmo & Fitness(2015--2024) & Fitness(2025) & Std & Evaluaciones & Tiempo (s) \\
\hline
ILS-ES      & -3.311 & -3.991 & 0.038 & 2863 & 0.023 \\
ILS-ES+CHC  & -3.297 & -3.984 & 0.028 & 6507 & 0.056 \\
\hline
\end{tabular}
}
\end{table}

\subsubsection*{S\&P 500 ($n=457$)}

\begin{table}[H]
\centering
\caption{Comparación ILS-ES vs.\ ILS-ES+CHC -- S\&P 500}
\label{tab:results_ilses_chc_sp500}
\resizebox{\textwidth}{!}{%
\begin{tabular}{lccccc}
\hline
Algoritmo & Fitness(2015--2024) & Fitness(2025) & Std & Evaluaciones & Tiempo (s) \\
\hline
ILS-ES      & -4.035 & -4.761 & 0.092 & 546 & 0.026 \\
ILS-ES+CHC  & -4.007 & -4.711 & 0.093 & 769 & 0.039 \\
\hline
\end{tabular}
}
\end{table}

\subsection{Análisis de Dispersión (Boxplots)}

\begin{figure}[H]
  \centering
  \includegraphics[width=\linewidth]{\detokenize{figuras/boxplot_comparativo_ilses_chc.png}}
  \caption{Distribución de fitness en entrenamiento (2015--2024) para ILS-ES (izquierda) e
  ILS-ES+CHC (derecha) en los tres mercados.}
  \label{fig:boxplot_ilses_chc}
\end{figure}

\subsection{Análisis de Resultados}

\subsubsection*{IBEX 35 ($n=30$): diversidad sin ventaja de fitness}

En el mercado más pequeño, el mecanismo de catástrofe CHC logra reducir la desviación estándar de
forma notable (de 0.014 a 0.008, un 47\% de reducción), reflejando mayor consistencia entre
ejecuciones. Sin embargo, el fitness de validación en 2025 es marginalmente peor (-3.129 frente a
-3.127 de ILS-ES). Este comportamiento es coherente con la teoría: en espacios de baja dimensión
con pocos activos, el ILS-ES ya converge con suficiente rapidez hacia óptimos de buena calidad, y
la macromutación del 40\% puede destruir estructuras parcialmente buenas sin compensar en el
presupuesto disponible. El coste en evaluaciones (2.7$\times$ más) tampoco se justifica aquí.

\subsubsection*{S\&P 100 ($n=97$): punto de equilibrio óptimo}

En el mercado de dimensión intermedia, ILS-ES+CHC mejora simultáneamente el fitness de
entrenamiento (+0.014) y el de validación (+0.007), al tiempo que reduce la variabilidad (Std
de 0.038 a 0.028, un 27\%). Este es el escenario donde el mecanismo de divergencia resulta más
efectivo: el espacio de búsqueda es suficientemente grande como para que el ILS-ES estándar quede
atrapado en cuencas subóptimas, pero la macromutación consigue redirigir la trayectoria hacia
regiones genuinamente distintas antes de que se agote el presupuesto.

\subsubsection*{S\&P 500 ($n=457$): mayor beneficio relativo en alta dimensión}

En el problema de mayor dimensión, ILS-ES+CHC obtiene su mejor mejora relativa: +0.028 en
entrenamiento y +0.050 en validación (un 1.06\% en valor absoluto). Resulta especialmente
significativo que este beneficio se consiga con un incremento moderado de evaluaciones (solo el
41\% más: de 546 a 769), lo que indica una alta eficiencia de la divergencia por catástrofe en
espacios de alta dimensionalidad. La desviación estándar permanece prácticamente estable (0.092
frente a 0.093), sugiriendo que la macromutación no introduce inestabilidad adicional.

\subsubsection*{Interpretación del mecanismo de divergencia}

Los resultados muestran un patrón claro que crece con la dimensión del problema. Definiendo la
eficiencia de convergencia como la mejora de fitness en validación por evaluación adicional empleada:

\begin{equation}
\text{Eficiencia} = \frac{\Delta f_{\text{2025}}}{\Delta \text{Evals}} =
\frac{f_{\text{CHC}} - f_{\text{ILS-ES}}}{\text{Evals}_{\text{CHC}} - \text{Evals}_{\text{ILS-ES}}}
\end{equation}

\begin{center}
\begin{tabular}{lrr}
\hline
Mercado  & $\Delta f_{\text{2025}}$ & Eficiencia \\
\hline
IBEX 35  & $-0.002$ & $-4.0 \times 10^{-7}$ \\
S\&P 100 & $+0.007$ & $+1.9 \times 10^{-6}$ \\
S\&P 500 & $+0.050$ & $+2.3 \times 10^{-4}$ \\
\hline
\end{tabular}
\end{center}

La eficiencia crece varios órdenes de magnitud al aumentar la dimensión, confirmando que el
mecanismo CHC está especialmente indicado en problemas de alta dimensionalidad donde el ILS-ES
estándar tiende a la convergencia prematura.

Otro hallazgo relevante es la \textbf{reducción consistente de la desviación estándar} en IBEX 35
y S\&P 100, que sugiere que el mecanismo de divergencia actúa como un \emph{piso de calidad}:
incluso en las ejecuciones menos afortunadas, el reinicio forzado rescata al algoritmo de los
óptimos locales más pobres.

\subsubsection*{Conclusiones}

\begin{enumerate}[noitemsep]
  \item \textbf{Beneficio creciente con la dimensión:} el mecanismo CHC es más eficaz conforme
  aumenta el número de activos. Para $n < 50$, ILS-ES estándar resulta más económico sin pérdida
  significativa de calidad; para $n > 100$, ILS-ES+CHC es recomendable.
  \item \textbf{Mayor robustez:} la desviación estándar se reduce en todos los mercados excepto
  S\&P 500, donde permanece estable. El algoritmo es más predecible en la práctica.
  \item \textbf{Equilibrio exploración--explotación mejorado:} el contador de estancamiento permite
  al algoritmo alternar de forma adaptativa entre intensificación (ILS normal) y diversificación
  (catástrofe CHC), en lugar de comprometerse permanentemente con una sola estrategia.
  \item \textbf{Coste computacional razonable:} el incremento en evaluaciones es de 2--3$\times$ en
  dimensiones bajas y solo del 41\% en S\&P 500, donde el beneficio es mayor.
\end{enumerate}
